\begin{abstract}\phantomsection\label{sec:resumen}
\noindent Sea \(Q(n)\) el número de formas de colocar \(n\) reinas sin ataques mutuos en un tablero
de \(n\times n\). El predictor canónico \(Q_{\rm pred}(n)\), una fórmula finita construida a partir del
maximizador de una entropía en el tablero, reproduce el término principal y la
constante de \(\log Q(n)\); ambos hechos proceden de preprints anteriores del autor. Determinamos ahora la
primera corrección al predictor finito:
el error con signo \(R_Q(n)=\log\bigl(Q(n)/Q_{\rm pred}(n)\bigr)\) satisface
\(nR_Q(n)\to\cerr\). Damos una fórmula explícita para este límite y obtenemos
\(Q(n)/Q_{\rm pred}(n)=1+\cerr/n+o(1/n)\).

\smallskip\noindent
Para obtener la fórmula elegimos una solución uniformemente al azar y revelamos sus filas en orden
aleatorio. El límite reúne una corrección inicial, la integral del defecto medio que deja cada
inserción y tres contribuciones del tablero que queda al detener el proceso: la de las torres, la de
un par de casillas en una misma diagonal y la de la interacción de las dos diagonales de un
cuadrado. Comparar dos niveles de parada da el límite y prueba que la integral converge
absolutamente. La demostración se apoya en resultados de esos preprints y está asistida por ordenador en
dos pasos.

\smallskip\noindent
La evaluación numérica de la fórmula da \(\cerr\approx0{,}40550038\), con una estimación empírica del
error de \(3\cdot10^{-8}\) y una banda conservadora \(0{,}4055004\pm2\cdot10^{-7}\). Son estimaciones
numéricas y no intervalos certificados: el signo y los decimales de \(\cerr\) siguen pendientes de
certificación.

\medskip\noindent
\textbf{Palabras clave:} problema de las $n$ reinas; enumeración asintótica; permanentes;
problema de torres; aritmética de intervalos.

\noindent
\textbf{MSC 2020:} 05A16; 60C05; 15A15.
\end{abstract}
